% 1. Thêm giải thích các thành phần cho công thức (2.1)
% 2. Sửa phần giải thích các thành phần công thức (2.2) và (2.3) rõ ràng hơn, đặt dưới công thức để dễ theo dõi


% Chỉ bổ sung thêm và không sửa cấu trúc của báo cáo
% Sửa "tẩng" thành "cấp"
% Mục tiêu c1 chưa đầy đủ
% Phát biểu bài toán: 
% phát biểu 2 cấp đều là tuyến tính chấp nhận giảm bớt độ chĩnh xác;
% Bổ sung:Cách tiếp cận baseline là quy hoạch tuyến tính quy hoạch tuyến tính có thể giải 2 cấp bằng cách đưa về KKT (Bổ sung vào cả chương 3 và chương 4)



% Code: Tạo và run ≈ 5000 kịch bản thuộc các nhóm khác nhau để đánh giá cả baseline quy hoạch tuyến tính và cách tiếp cận bi-level trên các kịch bản khác nhau. Quy hoạch tuyến tính có thể giải 2 cấp bằng cách đưa về KKT. Cần giải thích rõ hơn về cách tạo kịch bản và cách đánh giá hiệu quả của các phương pháp trên các kịch bản đó.
% yêu cầu: (viết chi tiết hơn)

13/6/2026:
Cập nhật hình 3.1, lấy ảnh D:\BS_thesis\DATN\Report\bao-cao\images\pipeline.png
Các phương pháp Bilevel-BO, PPO-500 đề xuất nếu kết quả tốt thì in đậm và gạch chân trong phần bảng kết quả

Các hình ở chương 5 đang bị lỗi font chữ, kiểm tra và sửa lại

KKT-LP đóng mở ngoặc baseline, các phương pháp khác đóng mở ngoặc đề xuất

Kết luận và tài liệu tham khảo để chương riêng như mở đầu. (kiểm tra và sửa lại)

Hạn chế nhất có thể cách viết sử dụng Bullet point và việc chêm xen tiếng anh trong câu